% ============================================================
% 6. LIMITATIONS -- humanized. Strong section already; just vary cadence,
% cut em-dashes, let some sentences be plain.
% ============================================================
\section{Limitations}

We put the limitations before the discussion's conclusions travel too far,
because several of them set exactly how far those conclusions can go.

\paragraph{Single-participant evaluation.}
The binding limitation is scope. The EK100 probe was evaluated on P01 alone
(\num{870} validation clips), and the HD-EPIC direct probes use P01 as well.
\S\ref{subsec:protocol} explains that this followed from which videos were on disk
at run time rather than from a deliberate choice. Every anticipation claim in the paper is therefore a claim about one
participant. Two analyses are exceptions and should be read as such: the
training-side ablations of \S\ref{subsec:training-ablation} and the
recoverability probes of \S\ref{subsec:horizon} train on P01--P07 and evaluate
on held-out participants, so where our conclusion rests on those it is not
single-participant. We make no assertion about
cross-participant generalisation, and the small spread between seeds should be
read as a single-participant artifact, not as evidence of stability across the
dataset.

\paragraph{Sample size below the generalisation threshold.}
Participant count aside, the direct probe runs on about \num{1900} training
examples against a high-dimensional input, and both arms interpolate that
training set (\S\ref{subsec:ddn}). In that regime held-out performance is set by
sample size, not by model or optimiser choices. This is why we call the result a
faint signal under a sample ceiling rather than no signal at all: the experiment
as run cannot separate a genuinely absent effect from a real one too small to
show at this $n$.

\paragraph{Mechanism claims rest on the direct probe, not the benchmark.}
The dissociation and the signal-content diagnostics come from the signal-present
HD-EPIC venue. The EK100 result is signal-absent by construction, since the
behavioral slots are masked at evaluation, so it speaks to whether a
behaviorally-trained predictor transfers, not to the gaze/hand-to-motion
mechanism. We do not carry evidence across that line. The cost of that discipline
is that our mechanism claims are never checked on the benchmark task itself.

\paragraph{The vision comparator sits below its own prior.}
Even fitted fairly (\S\ref{subsec:dissociation}), the vision arm does not clear
the verb class prior non-degenerately, and the behavioral arm never clears it at
all. The dissociation is established below the level of useful anticipation. A
stronger comparator, or a task with more headroom above the prior, might change
its size, though on this evidence not its sign.

\paragraph{The stronger pathway was not run.}
We retired the GazeQwen-informed pathway at the design stage
(\S\ref{subsec:b4}) on the strength of the diagnostics, without implementing it.
The diagnostics argue convincingly that it would not recover the signal, but the
argument is inferential. We did not falsify it directly, and a zero-initialised
gated implementation is still the clean way to do so if the sample-size
constraint is ever lifted.

\paragraph{The signal-drop control is scored on the conditioning objective, not on EK100.}
The matched never-conditioned control of \S\ref{subsec:training-ablation} is a
full-rate ablation of the very checkpoint that carries every anticipation result
here, which is stronger than a comparison across models. It is scored, however, on
feature-space prediction error over a held-out participant rather than through the
EK100 probe itself. What remains unrun is the same ablation carried all the way
through to anticipation recall; we therefore report the null in signal content as
established on the conditioning objective and inherited, rather than
independently re-measured, downstream.

\paragraph{The checkpoint is under-trained relative to the intended scale.}
Every anticipation number here rests on a predictor fine-tuned for three epochs at
thirty clips per recording (\S\ref{subsec:method-setup}). That is the small
configuration, not the scale the conditioning study was designed around, and no
full-scale run ever completed. This matters in a specific direction: a null
measured on an under-trained predictor is weaker evidence than it reads as,
because a model that has not finished learning the task has not finished learning
what to do with an auxiliary input either. Two things limit how much weight this
objection can carry. The feasibility curve of \S\ref{subsec:protocol} is
essentially flat after the fifth epoch, and the fine-tune is demonstrably worth
$0.125$ MSE over the stock predictor (\S\ref{subsec:training-ablation}), so the
training did take. But we cannot rule out that a longer schedule would give the
conditioning pathway something it currently lacks, and we do not claim otherwise.

\paragraph{Goal-level probing is infeasible on one participant.}
The representation-level test we wanted to run, a recipe-goal probe on
HD-EPIC~P01, does not work at single-participant scale. Under the hand-validity
session split, four of eight recipes have no held-out examples, one recipe
accounts for about three-quarters of the held-out set, and recipe identity is
nearly collinear with session identity. That is a characterised data requirement
rather than a failed experiment, and it is what motivates the multi-participant
direction of \S\ref{sec:conclusion}.